\chapter{Formal Safeguards with TLA+}
\label{chap:formal-safeguards}
\newcommand{\tlbreak}[2]{\texttt{#1}\allowbreak\texttt{#2}}

CharlotteOS uses executable TLA+ models to review stateful functionality. The
models state each safety contract, explore action orderings in a finite instance,
retain known-bad designs as counterexamples, and map abstract transitions back to
Rust.

This is stronger than relying on selected test schedules, but narrower than a
proof of the complete operating system. TLC proves invariants of the modeled
state machine for the configured finite bounds. The repository supplements
that result with an action-level conformance map, Rust regression tests,
target builds, Clippy, and boot tests. Each layer answers a different question:

\begin{description}
    \item[TLA+ model] Is the protocol design safe under every modeled action
        ordering, crash point, retry, and stale event within the configured
        bounds?
    \item[Conformance map] Which concrete Rust state and linearization point
        implement each abstract action?
    \item[Rust and boot tests] Does the implementation exercise the intended
        transition on the host and on CharlotteOS targets?
    \item[Build and lint gates] Does the integrated implementation remain
        type-correct and warning-free for the actual target?
\end{description}

The models and their configurations live in \repodir|docs/tla|. The concise
model inventory is in \repofile|docs/tla/README.md|; the concrete-to-abstract
mapping is in \repofile|docs/tla/CONFORMANCE.md|. Those files are maintained
with the implementation and are authoritative over the overview in this
chapter.

\section{Why model functionality as a state machine?}

Most failures addressed by these models are not computation errors. They are
ordering errors: a stale cleanup arrives after a replacement, a thread is
removed while it can still execute, a session identifier is reused after a
restart, or a joiner forgets which cluster has authority over it. A unit test
can exercise one such interleaving. A state-machine model asks TLC to enumerate
all reachable interleavings in a deliberately small world.

A TLA+ safety specification has four essential parts:

\begin{enumerate}
    \item \textbf{Variables} record only the state relevant to the property,
        such as an owner, generation, durable term, queue, or admission anchor.
    \item \textbf{Init} defines every allowed initial state.
    \item \textbf{Next} is a disjunction of named actions. Each action defines
        its enabling condition and complete next-state update.
    \item \textbf{Invariants} state what must be true in every reachable state.
\end{enumerate}

The common shape is:

\begin{verbatim}
VARIABLES phase, generation, owner
vars == <<phase, generation, owner>>

Init == ...
Next == Start \/ Replace \/ Stop \/ Crash \/ Restart
Spec == Init /\ [][Next]_vars

TypeOK == ...
Safety == ...
Invariants == TypeOK /\ Safety
\end{verbatim}

The square temporal operator means that every step must be a \texttt{Next}
step or a stuttering step and that the invariant must hold throughout the
behavior. TLC constructs the reachable state graph rather than sampling it.
Constants in a \texttt{.cfg} file bound identities, queue capacities,
generations, terms, and log indices so that this graph is finite.

Small bounds are intentional. Many protocol violations need only two owners,
two generations, or three Raft nodes. Exhaustive exploration of that small
instance is often more revealing than a large randomized run because TLC
systematically covers the orderings rather than hoping to encounter them.

\section{The safeguarding loop}

CharlotteOS applies the models as a feedback loop between design,
implementation, and validation:

\begin{enumerate}
    \item \textbf{Choose a safety boundary.} Identify the authority or resource
        whose invalid transition would be harmful: reply authority, memory
        ownership, an executing address space, a committed log entry, or a
        service generation.
    \item \textbf{Project concrete state.} Remove bytes, addresses, and policy
        details that cannot affect the selected invariant. Preserve identity,
        ownership, durability, ordering, and lifecycle phases.
    \item \textbf{Name the transitions.} Actions use implementation language.
        Names such as \texttt{Receive}, \texttt{Reap}, and \texttt{Restart}
        make review against Rust possible.
    \item \textbf{State the invariant positively.} Examples are ``a reaped
        thread is off-CPU,'' ``a committed entry agrees at every replica,'' and
        ``a stale generation cannot remove its replacement.''
    \item \textbf{Model crashes and stale work.} Durable and volatile state are
        separated where restart matters. Delayed messages, old handles, and
        retries are actions rather than assumptions that the environment is
        well behaved.
    \item \textbf{Retain a negative specification.} When a faulty transition is
        understood, an \texttt{Unsafe...} action and configuration are kept.
        The check suite requires TLC to find the named invariant violation.
        This proves that the invariant and test harness are capable of seeing
        the regression rather than passing vacuously.
    \item \textbf{Map back to code.} \repofile|docs/tla/CONFORMANCE.md| records
        the concrete function, persisted field, and linearization point for
        every modeled action. Disagreement is resolved by changing the model,
        changing CharlotteOS, or documenting behavior outside the abstraction.
    \item \textbf{Validate both levels.} Run TLC, focused Rust tests, target
        builds and Clippy, and relevant boot or multi-node tests.
\end{enumerate}

We use the executable model to check the safety contract alongside the
implementation. A change is incomplete if
the model claims an atomic or durable transition that the code does not
provide.

\section{What is modeled}

The current suite checks thirty-six safe configurations across thirty-four
specifications. It also checks 107 expected-counterexample
configurations. The models are divided by state ownership and failure boundary
so that each remains reviewable.
The Raft models are deliberately layered: election, log replication,
membership, pre-membership admission, and snapshots have different state and
different proof obligations.

\small
\begin{longtable}{|p{0.26\textwidth}|p{0.26\textwidth}|p{0.38\textwidth}|}
\hline
\textbf{Model} & \textbf{Functional boundary} & \textbf{Principal safety properties} \\
\hline
\endfirsthead
\hline
\textbf{Model} & \textbf{Functional boundary} & \textbf{Principal safety properties} \\
\hline
\endhead
\tlbreak{Charlotte}{IPC} & Endpoint calls, reply tokens, memory move/copy/borrow,
    cancellation, endpoint close, domain teardown & Bounded queues; unique
    receiver-visible reply authority; exclusive write borrow; every remaining
    borrow backed by a live token; no dangling moved-memory capability. \\
\hline
\tlbreak{CharlotteIpc}{Transaction} & Multi-entry vector preparation,
    connection attachment, rollback, delivery, reply and non-terminal timeout &
    Failed preparation restores every affected registry; attached connections
    cannot escape their transaction; timeout cannot release a live call's loan. \\
\hline
\tlbreak{CharlotteIpc}{Publication} & Captured reservation, move escrow,
    hidden memory/connection delivery and namespace replacement & Publication
    validates the whole batch; receive/observation grants access once; late
    cleanup preserves successor authority and original metadata charges. \\
\hline
\tlbreak{Charlotte}{EndpointObservers} & Readiness observers and connection-close
    watches & Ordinary message arrival cannot falsely signal endpoint closure;
    closure wakes the appropriate observers. \\
\hline
\tlbreak{Charlotte}{CQ} & Completion operations, CQ ring/backlog, wait/wake,
    cancellation, timers & Non-lossy completion tracking; bounded visible ring;
    no lost wakeup; timer completion remains tied to its operation; a buffered
    operation retains its kernel loan through cancellation until terminal completion. \\
\hline
\tlbreak{Charlotte}{TimedWait} & Work-generation publication, delayed wake
    delivery, and watchdog expiry & A timer cannot report timeout after a new
    work generation has already been published. \\
\hline
\tlbreak{Charlotte}{CheckedWait} & Fallible watchdog/source admission,
    registration, park recheck, unrelated wakes and result observation &
    Pending owners and loans survive rejected admission and timeout; wakes
    cannot adopt results; retries preserve one deadline and accepted request. \\
\hline
\tlbreak{Charlotte}{Scheduler} & Spawn, admission, execution, blocking, migration,
    cross-LP abort, reaping, address-space destruction & At most one execution
    per LP; only the owner LP retires an executing thread; reaping occurs
    off-CPU; domain abort prevents late sibling publication; destroyed address
    spaces have no live threads. \\
\hline
\tlbreak{Charlotte}{ThreadJoin} & Recyclable TIDs, spawn generations, delayed
    exit-observer registration & A captured handle observes only its exact
    generation; a missing or recycled slot completes immediately; exhausted
    generation identity fails closed before wrap. \\
\hline
\tlbreak{Charlotte}{StackGrowth} & Demand-grown user stacks, admitted budgets,
    pressure reserve and domain exit & Growth stays within its committed-page
    budget and preserves the free-frame reserve; modeled exit releases frames. \\
\hline
\tlbreak{Charlotte}{AddressSpace} & Reusable numeric ASIDs and software-generation
    handles & A stale handle cannot close a replacement that reused the same
    numeric ASID. \\
\hline
\tlbreak{Charlotte}{RootRetirement} & Exact-root leases, staged close,
    invalidation retry, abandonment and physical release & Complete owners
    survive uncertain invalidation; leases fence detachment; physical cleanup
    runs once after acknowledgement; rejected release retains original charges. \\
\hline
\tlbreak{Charlotte}{HardwareAsid} & Hardware-ASID allocation, retirement, cached
    translations, invalidation & A hardware tag is not reused while stale
    translations for its former owner remain. \\
\hline
\tlbreak{Charlotte}{InterruptRoute} & Interrupt binding, queued wake, unbinding,
    deferred drain & A queued wake cannot be delivered through a later route
    generation to the wrong owner. \\
\hline
\tlbreak{Charlotte}{Capability} & Unified per-address-space capability namespace
    and move rollback & Handles are fresh and authoritative; move revokes the
    source; rollback restores only the exact transaction; a closed namespace
    is empty. \\
\hline
\tlbreak{Charlotte}{Authorization} & Versioned subject/service policy,
    generation-fenced service bindings, decisions and capability issuance &
    Only authorized roles mutate policy or publication; decisions are bound to
    a subject, policy version, service generation and rights; redemption is
    single-use and cannot amplify authority. \\
\hline
\tlbreak{CharlotteName}{Publication} & Local public/policy classes, prepared
    death watches, replacement and retirement & Public paths cannot issue
    policy-published connections; watch admission precedes mutation; stale hints
    need the current producer's own result; disposal preserves watch/source order. \\
\hline
\tlbreak{CharlotteBoot}{Trust} & Verifier, exact installed fixture, one-shot
    handoff and conditional protected installation & Signed verification does
    not install policy; boot uses committed identity; the slot cannot be replaced;
    future commit/reboot ordering assumes protected atomic state. \\
\hline
\tlbreak{Charlotte}{DMA} & CPU mapping, IPC loans, coherent/exclusive SMMU
    mapping, invalidation, quarantine, driver exit, pin reclamation & Unique
    stream/domain ownership; exclusive DMA has no overlapping CPU authority,
    loan, or second pin; mapped memory stays pinned; failed hardware
    acknowledgement retains rather than prematurely frees authority. \\
\hline
\tlbreak{CharlotteDma}{Operation} & Fixed-unit readiness, retained map/unmap
    owners, typed rejection and CPU teardown & Availability grants no completion
    proof; started attempts cannot replay; confirmed DMA survives CPU-unmap
    failure; uncertain pins retain backing and its original charge. \\
\hline
\tlbreak{Charlotte}{ServiceLifecycle} & Signed artifact staging, loading, start,
    two-phase publication, lookup, fenced unregister, exit and teardown &
    Untrusted images never load; only activated generations resolve; stale
    activation/unregister cannot affect a replacement; teardown follows
    domain-wide reaping after cooperative exit or domain abort. \\
\hline
\tlbreak{Charlotte}{Raft} & Terms, durable votes, election, higher-term step-down,
    crash/restart & A voter records at most one durable vote per term and at
    most one leader can hold a term. \\
\hline
\tlbreak{Charlotte}{RaftLog} & Append, conflict repair, replication, commit,
    crash/restart & Log matching, committed-entry agreement, commit within the
    log, and leader completeness. \\
\hline
\tlbreak{Charlotte}{RaftResponse} & Delayed term replies, admitted peers and
    confirmed replication progress & An old prefix acknowledgement cannot commit
    a replacement; reordered replies and backoff preserve confirmed progress. \\
\hline
\tlbreak{Charlotte}{RaftRead} & Per-read log fences, delayed replies, ownership
    binding and query completion & Reads require a fresh quorum-confirmed index,
    application and exact node/domain, term and membership episode. \\
\hline
\tlbreak{Charlotte}{RaftMembership} & Voters/learners, joint consensus,
    finalization, decommissioning & Both voter majorities are required while
    joint; all proposed peers reach the joint fence before finalization;
    leadership and decommissioning follow active membership. \\
\hline
\tlbreak{Charlotte}{RaftJoin} & Selected-anchor admission before membership,
    JOIN fence, catch-up, crash/restart & A joiner accepts only the chosen
    anchor, cannot campaign, is promoted only after catch-up, and restores its
    durable admission fence after restart. \\
\hline
\tlbreak{Charlotte}{RaftSnapshot} & Chunk reception, durable installation,
    matching suffix, membership recovery & Stale snapshots cannot regress
    progress; the installed state and retained suffix agree; restart restores
    membership and application state before use. \\
\hline
\tlbreak{Charlotte}{RaftStream} & Leader/term-bound snapshot chunks and offset
    order & A replacement election discards partial bytes; installation uses
    one sender episode in strict offset order. \\
\hline
\tlbreak{Charlotte}{SnapshotRestore} & Application acceptance, rejected restore,
    persistence failure and restart exposure & Rejection preserves application,
    durable state, progress, membership and join admission; recovery must accept
    the snapshot before exposing its node. \\
\hline
\tlbreak{CharlotteCluster}{Ingress} & Joint membership, replica placement,
    exact-generation readiness, drain, immutable ingress snapshots, and
    retained flow epochs & New flows require a fresh snapshot lease; eligible backends are
    placed, ready, admitted, and non-draining; incomplete snapshots do not
    install; an unretained flow epoch cannot silently remap. \\
\hline
\tlbreak{Charlotte}{PlacementControl} & Filtered pressure, dwell, cooldown,
    eligibility loss and capacity reservations & Pressure-only moves respect
    stability gates; lost eligibility permits forced movement; reservations
    preserve the protected capacity reserve. \\
\hline
\tlbreak{Charlotte}{RemoteCall} & Call identity, target generation, duplicate
    suppression, timeout uncertainty, result eviction & At most one execution
    while tracked; stale targets never execute; uncertain outcomes are not
    reported as success; deduplication evidence remains until transport
    settlement. \\
\hline
\tlbreak{Charlotte}{ReliableMessage} & Wire sessions across retry abandonment and
    unilateral service restart & Logical epochs have unique ordered wire
    identities and delayed sessions cannot move receive state backwards. \\
\hline
\end{longtable}
\normalsize

This decomposition also states what is \emph{not} silently being claimed. For
example, Ethernet fragmentation and retry timing are not part of the Raft
election proof; protobuf bytes and peer addresses are not part of joint
consensus; cryptographic correctness is not part of the service-lifecycle
trust-bit abstraction. Such omissions are recorded in the conformance map.

\section[Model-driven findings]{How model reasoning changes CharlotteOS}

The most valuable result is often a mismatch discovered while constructing the
model. The correct response depends on which side was wrong: sometimes the
manual or model overstated functionality; sometimes the code lacked a safety
fence.

\subsection{Authorization is controlled capability issuance}

The local name service implements policy-controlled authorization alongside
public bootstrap naming. Bearer-key lookup opcodes are retired. The
signed loader binds an artifact-derived stable principal, exact address-space
generation, and roles in a kernel-only table. IPC enqueue places that
authenticated snapshot in the receive envelope; request memory cannot choose
the caller identity.

The authorization model makes the intended boundary explicit. The catalog
publishes a service generation and the maximum rights that may be delegated.
A separate logical policy role maps a stable principal and service identity to
an allowed-rights set and policy version. A decision is bound to the
kernel-authenticated subject, requested rights, current service generation and
policy version, and can be redeemed once for an attenuated connection. A
co-located name/policy process performs decision and attenuated connection
delegation as one serialized transition. Policy mutation and authorized
publication require distinct roles, and issuance, denial, and delegation
failure enter a bounded administrator-readable audit stream. A later split
service needs an unforgeable single-use grant and must revalidate every binding
at redemption. Policy and audit state are not yet durable or replicated.
Public lookup cannot resolve a policy-published service, and public
registration cannot displace one. These publication-class checks remain
outside the authorization model and are projected separately below.

The model also prevents an attractive but false claim. A policy update blocks
new issuance and invalidates unredeemed decisions, but cannot selectively
retract a direct connection already delegated by the kernel. Hard revocation
requires endpoint closure, a revocable proxy, or a future kernel lease object.
This prospective revocation contract is safer than documenting authority that
the implementation cannot enforce. The practical design is recorded in
\repofile|docs/architecture/authorization-policy.md|.

\subsection{Boot verification and installation are separate}

\tlbreak{CharlotteBoot}{Trust} distinguishes a signed verification receipt from
boot authority. The shared verifier accepts either the exact installed record
or its direct successor with the accepted digest as predecessor. The x86 boot
fixture requires the exact compiled installed revision and digest before
preparing its recipient-key owner. One-shot kernel publication then supplies
the public policy used by both service manifests and the operational gate.
A duplicate publication returns its unused preparation; development trust is
selected explicitly and is never a fallback after verification failure.

The model has two configurations. The first projects those implemented fixture
and development paths. The second is a conditional contract for the missing
protected installer: it assumes an authenticated context, atomic rollback-resistant
state commit and serialization against boot publication. Two signed successors
may verify against one old expectation; only one can commit it. An interrupted
commit may leave the whole old or new record durable, so the caller must read
protected state again before publishing. Recovery must first settle or fence the
original writer; reading old bytes alone proves no completion. Reboot preserves that committed state
while discarding volatile verification and publication state.

These assumptions describe work still needed for production. The model adds
no installer, firmware authentication, persistent-state adapter, key custody or
production selector. Its ten negative controls expose the weakened ordering,
identity and handoff guards. Byte-level cryptography, actual storage atomicity,
erasure, outer contexts and physical lifetime remain outside the projection.

\subsection{Local publication owns its death watch}

\tlbreak{CharlotteName}{Publication} follows the node-local registry. A candidate
owns its source connection before preparing a persistent endpoint-death watch.
Watch rejection leaves the current publication authoritative. Successful policy
publication and registry replacement are separate internal steps; the serialized
reactor admits no request between them. Deferred replies use the installed
publication, and only then is the previous watch disposed before its source.
Public lookup rejects policy-published entries whether immediate or deferred.
Public replacement requires an existing public entry from the same nonzero
registrar. Authorized lookup still needs its policy decision; both publication
classes have policy-store bindings.

A CQ entry supplies a readiness hint. Even when an old hint matches a current
watch, retirement needs that current owner's own confirmed endpoint-close
result. Pending or unrelated hints preserve the entry. Unexpected observation,
rejected waiting and lifecycle abort retain the registry through terminal abort.
Explicit retirement separately requires the exact generation and authenticated
registrar or manager. These transitions establish logical unpublication and
supply no physical cleanup proof.

Source comparison found a further exhaustion boundary. The policy store could
publish an unsigned generation beyond the positive signed NS result range.
The subsequent NS conversion rejected after policy mutation. Publication now
checks the signed limit first, preserving an active binding or tombstone on
exhaustion. Host regressions reproduced the earlier behavior and verify that
\texttt{i64::MAX} can be published once before later attempts reject.

One finite instance checks two names, principals and candidates at live capacity
one; a smaller three-candidate projection reaches generation exhaustion at two.
Thirteen required counterexamples weaken the publication, retirement and ownership
guards. Kernel authentication and watch-result provenance, trusted launch-budget
preparation, wire parsing, allocation/destructor contexts and physical completion
remain boundary assumptions.

\subsection{IPC attachments are cross-registry transactions}

The base IPC model intentionally collapses vector attachments. The
\tlbreak{CharlotteIpc}{Transaction} refinement instead represents two or more
moved entries, a live memory loan, an attached connection, the queue and the
pending-call registry. Current vector calls carry the memory entries, whereas
scalar calls carry connection attachments and may copy memory. The refinement
deliberately composes these paths into a stronger rollback obligation; it does
not claim that one current syscall carries every form. No request is visible
until preparation commits, and every pre-commit failure restores all affected
registries. Its unsafe rollback keeps one moved entry and a connection in the
target after failure and must violate \texttt{RollbackRestoresAll}.

A reply-wait timeout is also modeled separately from terminal call completion.
It reports only that no reply was observed during the wait and therefore cannot
release a buffer still reachable by the kernel or receiver. The unsafe timeout
transition releases that loan while the call remains pending and must violate
\texttt{TimeoutRetainsLoan}.

\tlbreak{CharlotteIpc}{Publication} follows the owning outputs through a
further boundary. A unified capability entry can be live while its memory or
connection payload remains hidden in a queued request or unobserved reply.
Receive commit and first reply observation make those resources usable.
Readiness alone changes no ownership. Failed receive admission or result
writing preserves the queued attachments, and repeated observation cannot
revive a returned resource that its caller has already closed.

The same model retains captured namespace accounts through reservation,
source escrow and detached metadata disposal. Rollback restores the existing
source slot without fresh admission, including at its record ceiling. A
replacement may reuse the numeric ASID and handle values; late cleanup must
still affect only its original account. Physical cleanup proofs, loans and
actual allocator/IRQ/guard contexts remain separate obligations.

\subsection{DMA availability and completion}

\tlbreak{CharlotteDma}{Operation} separates unit availability, backend
acknowledgement and public completion. Readiness retains the original map
owner until both the actual command engine and recorded domain cell can be
claimed. Preparation may lose a race to restoration, and another claimant may
win before a wake is processed. The owner rechecks the same unit and starts
mapping work once. Busy waiting requires a qualified scheduled caller;
masked ready paths can still complete without waiting.

Before entering unmap, the runtime records the attempt as retained. Only a
typed rejection before backend claim can rearm it. A failure after work starts,
including a busy-looking final completion error, grants no replay permission.
Confirmed DMA completion survives failed CPU teardown, so a later finish can
retry CPU unmap without repeating DMA unmap. Failed shared setup keeps its CPU
mapping behind a restricted owner; local teardown cannot refund backing that
an uncertain DMA pin still retains. Hardware acknowledgements, physical release
and actual allocator/IRQ/guard contexts remain separate obligations.

Source comparison found a missing step in that abstraction. Consuming the
restricted wrapper through \texttt{unmap} returns an unmapped memory owner.
The kernel had fenced fresh CPU mappings only for exclusive DMA pins, so a
retained coherent pin could still allow a fresh mapping with safe Rust access.
A controlled kernel regression reproduced that path. Fresh read-only and
writable mappings now require all DMA pins to be gone. Existing coherent
mappings are installed before pinning and keep their volatile-access contract.
The model includes fresh remapping and a required counterexample that removes
this fence; the audit records the failing ARM run and passing ARM/Intel runs.

\subsection{Two-phase service publication}

An earlier lifecycle abstraction treated replacement publication as one
atomic swap. The implemented distributed catalog instead performs three
observable steps: commit an inactive prepared generation, publish the local
connection, and commit activation. Lookup intentionally excludes the inactive
generation, so replacement has a temporary unavailable interval.

The model was changed to describe that real contract rather than claim
continuous availability. It also exposed the code-level obligation that
generation allocation must fail closed. Distributed generation saturation and
local unchecked increment were replaced with checked allocation that returns
an error without mutating the catalog. A delayed unregister is accepted only
when both owner and generation still match, which preserves a replacement.

The model records this availability limit while preserving the generation
safety that the implementation enforces.

\subsection{Reliable-message session identity}

The reliable-message service formerly initialized a wire session from the
service generation and incremented that number when abandoning a retry epoch.
Generation $N$ after one abandonment could therefore use the same identity as
generation $N+1$ before abandonment. A bounded retired-session window also
allowed a sufficiently old delayed SYN to replace newer receive state after it
fell out of the window.

The session model represents a session as the ordered pair
$(\mathit{serviceGeneration},\mathit{retryEpoch})$. The repaired wire encoding
packs those values into disjoint 32-bit fields, accepts a replacement session
only when the packed identity is well formed and strictly newer, and disables
sends rather than wrapping at exhaustion. Two unsafe configurations retain the
old flat-identity collision and receive-regression behaviors as required
counterexamples.

\subsection{Restart-safe Raft admission}

Dynamic admission originally kept \texttt{joining} and
\texttt{joining\_from} only in process memory. A node could select an existing
cluster as its authority, crash before membership became durable, restart from
its launch-time singleton configuration, and campaign. The gap represented
a behavior of the implementation that the model had not captured.

The repaired implementation persists a \texttt{JoinAdmission} record
containing the selected anchor and the snapshot index that existed before
admission. Construction restores the joining posture and disables elections.
When the joint membership commits, the node first writes a newer
membership-bearing snapshot and only then clears the durable admission record.
This ordering covers both crash edges:

\begin{itemize}
    \item before the membership snapshot is durable, restart retains the
        anchor fence and the node cannot campaign;
    \item after the snapshot is durable but before the fence-clear write,
        restart recognizes the strictly newer membership snapshot as completion
        evidence and safely clears the stale admission record.
\end{itemize}

Operational auto-join belongs to the durable DNS-owned Raft node. Volatile
Raft nodes must not auto-join: without stable state, a process restart cannot
prove which cluster has authority. The safe
TLA+ specification includes crash and restart; an unsafe restart that forgets
admission must violate \tlbreak{RestartPreserves}{Admission}.

\subsection{Membership changes preserve the tested Graft contract}

Charlotte's Raft port derives from Graft, whose Java, C++ and Rust runtimes are
intended to form mixed clusters. The upstream Jepsen harness exercises hostile
schedules and membership changes. That evidence makes core and runtime parity
an explicit review concern; a local model or host test does not replace it.

A source comparison found that Charlotte had replaced the Java/Rust automatic
finalization catch-up gate with immediate proposal after JOINT committed. The
port now waits for every proposed remote voter and learner to reach the committed
JOINT fence. The leader owns that prefix locally. Committing FINALIZE still
requires both voter majorities. A slow proposed member can delay automatic
completion; explicit/manual finalization is a separate path.

A second regression finalizes the current leader as a retained learner. The
previous Charlotte code kept it in the leader state. It now becomes a follower
immediately, keeps its existing-term durable vote, and remains a member. The
upstream Jepsen demotion scenario selects a non-leader, while leader removal has
its own scenario; the new regression covers the retained-leader transition.

\tlbreak{Charlotte}{RaftMembership} checks 6,232 distinct states. It includes
next-only voter elections with both majorities and rejects automatic finalization
while a proposed member lags. Three counterexamples remove the catch-up,
leader-eligibility and dual-quorum guards. Durable vote persistence remains a
source-test obligation. Native Charlotte checks do not establish Graft TCP
interoperation, mixed implementation histories or read-lease correctness.

\subsection{A delayed acknowledgement belongs to its original term}

A source regression captured an honest term-one acknowledgement for index three,
then replaced the local prefix under term two. On re-election in term three,
a new no-op occupied index three. The old response incorrectly committed the
replacement, although the responding voter still held the old bytes.

The core now requires exact current term and a matching admitted source before
crediting an append or completed-snapshot response. Successful indices cannot
exceed the local tail. Reordered successes retain the strongest confirmed prefix;
append rejection cannot back off below it. Committed joiners can still replicate
before promotion, while voter quorum rules remain unchanged.

\tlbreak{Charlotte}{RaftResponse} checks this boundary in 49,395 distinct states
and retains four counterexamples. The source tests reproduce the unsafe commit
before correction. Transport authenticity, request correlation, removal epochs
and the timing assumptions of the existing quorum-contact read gate remain
outside this proof.

\subsection{Each read carries a fresh quorum fence}

A recent reply cannot prove that a leader still has authority. The host audit
reproduced this with three nodes: n2 and n3 committed value 2, then an honest
term-one reply reached isolated n1 and reopened its receive-time query gate.
n1 returned value 1. Current-term commitment alone did not prevent that read.

The query API now admits each request before appending a fresh no-op log entry.
Its \texttt{PendingRead} owner captures exact node/domain identity, term, index
and membership epoch. The reactor retains it while processing quorum traffic.
Readiness observes commitment and application without executing the query;
completion rechecks those conditions and consumes the owner once. Earlier
acknowledgements cannot confirm its new index. Compaction preserves the proof
through a newer committed current-term anchor.

DNS keeps the original monotonic deadline and typed reply destination with the
fence. It admits bounded queue storage before publication. Placement lookup
returns metadata; a local result carries no connection authority. Ordinary NS
lookup strips delegation rights, so DNS does not attempt to re-mint that grant
for a different caller. Invocation starts only after confirmed routing and
another deadline check, using the existing owned NS lookup/call phases on the
shared CQ.
Timeout or drop abandons the query while its no-op may still commit; it cannot
execute or replay the business call. Readiness-binding rejection aborts with
accepted work live. No extra per-read timer or blocking reactor wait is added.

\tlbreak{Charlotte}{RaftRead} explores 9,916 distinct states, including delayed
replies and external leader replacement unseen by the old node. Four controls
remove commitment, application, fresh-fence or owner-binding checks. The former
stale-read host schedule now rejects. A quorum may have confirmed a fence after
admission and replied before a later election; the read proof concerns that
confirmation within the request, rather than authority at delivery time.

The shared no-op wire grammar is unchanged. Each read adds a durable log entry
and quorum commit, with corresponding storage and latency costs. No performance,
Jepsen or mixed Charlotte/Graft execution claim is made. Bounded status/ingress
observations retain receive-time freshness separately; this is not general lease
qualification.

\subsection{Elapsed time comes from the kernel counter}

The reactors previously counted one 25 ms tick after waking even if seconds had
passed. They now convert a kernel-monotonic counter from one fixed epoch,
sampling after waits and before inbound or scheduled work. Fractional ticks
accumulate; invalid or regressing samples abort with pending work live. Timers
wake the service without manufacturing elapsed time. Clock conversion, resource
readiness delivery and physical cleanup remain outside the read model.

\subsection{Snapshot chunks belong to one leader and term}

Snapshot chunk assembly also needs a sending identity. The old receiver matched
only the snapshot boundary and offset: a new leader, or the same leader in a new
election term, could complete an earlier buffer. Host regressions reproduced
installation of a combined prefix and suffix. The buffer now captures the leader
and RPC term, and elections and higher-term observations discard it. A replacement
must begin at offset zero. \tlbreak{Charlotte}{RaftStream} checks this owner and
strict ordering in 400 distinct states, with two weakened controls. It assumes
one immutable byte image per sender/term/boundary; same-episode image replacement,
transport authentication and general allocation failure remain outside this model.

\subsection{Rejected application snapshots leave known ground intact}

A complete snapshot stream can still contain an unsupported application format
or a malformed late section. The catalog parses all sections into temporary
collections before replacing live state. A rejected restore leaves those live
collections intact. Raft propagates rejection before publishing the snapshot,
commit/applied indices, membership or join completion. Construction returns an
error before exposing the node or clearing durable admission.

\tlbreak{Charlotte}{SnapshotRestore} makes this boundary explicit alongside the
existing chunk and suffix model. Its finite configuration explores 168 distinct
states; seven weakened configurations must produce counterexamples. Separate
streamed/recovery coverage requires both completed joins and retained admission.
This corrects an earlier assignment-precedence omission that excluded join
clearing while broader activation coverage still passed. An accepted
newer membership snapshot can complete a join, while the original snapshot or
membership excluding this node retains its admission fence.

Restore succeeds before the durable write. If persistence fails, the Raft caller
terminates; a later caller must restore the whole durable image again. This
ordering does not provide an atomic transaction across application and disk.
The model assumes the existing whole-object storage contract and serialized
Raft execution. It does not prove decoder correctness, fallible allocation,
concurrent-reader isolation, endpoint withdrawal or sibling-thread termination.

\subsection{Cluster placement and DSR ingress}
\label{subsec:formal-cluster-ingress}

The cluster-ingress model begins after a Raft application command commits. It
then composes stable and joint membership, desired replica sets, deployment
generations, per-node readiness, drain intent, discovery routes, asynchronous
snapshot installation, VIP advertisement, and per-flow epoch retention. This
keeps Raft quorum mechanics in their existing models while checking the point
where replicated authority becomes packet-routing policy.

The safe specification derives every eligible backend as the intersection of
active membership, committed placement, readiness for the exact deployment
generation, and nodes not draining. A router installs only a complete snapshot
for which all member routes are known. New flows require the router's installed
policy to hold a bounded local freshness lease. The policy may lag the latest
global commit during asynchronous propagation; an abstract fair expiry action
prevents an unrenewed lease from persisting forever. Established flows retain
their selected policy while its snapshot remains present. After history
eviction, the binding remains pinned but its packets fail closed instead of
silently choosing a new backend. Concrete snapshot installation additionally
requires a quorum-fresh Raft leader or a follower with a recent successful
leader-log match, so a responsive but partitioned local DNS process cannot
renew the lease indefinitely.

Three negative configurations expose cross-layer failures. One
admits a new SYN through a stale snapshot after its freshness lease expires.
Another evicts an epoch still named by a live flow and then remaps that
flow through the current snapshot. The third admits readiness from an old
deployment generation. Exact readiness is already enforced by the catalog.
The router now enforces the other two contracts with a five-second monotonic
snapshot lease and fail-closed missing-epoch tombstones. All three unsafe
actions remain as regression witnesses.

\subsection{Retained regression patterns}

Other negative specifications preserve compact versions of repaired faults and
present conformance obligations:

\small
\begin{longtable}{|p{0.27\textwidth}|p{0.31\textwidth}|p{0.32\textwidth}|}
\hline
\textbf{Unsafe action} & \textbf{Failure represented} & \textbf{Required violation} \\
\hline
\endfirsthead
\hline
\textbf{Unsafe action} & \textbf{Failure represented} & \textbf{Required violation} \\
\hline
\endhead
\tlbreak{Unsafe}{MessageWake} & Message readiness and endpoint-close observation
    share one observer queue. & \tlbreak{CloseSignalImplies}{Closed} \\
\hline
\tlbreak{Unsafe}{Rollback} & A failed abstract attachment transaction leaves one
    moved entry and its attached connection in target registries. &
    \texttt{RollbackRestoresAll} \\
\hline
\tlbreak{UnsafeTimeout}{Release} & Reply timeout releases a loan while the
    pending call can still access it. & \texttt{TimeoutRetainsLoan} \\
\hline
\texttt{UnsafeUseMemory} & A guessed memory handle is used before delivery. &
    \texttt{UseRequiresDelivery} \\
\hline
\tlbreak{UnsafeUse}{Connection} & A guessed connection handle is used before
    delivery. & \texttt{UseRequiresDelivery} \\
\hline
\tlbreak{UnsafePublish}{Prefix} & One attachment becomes live before the whole
    batch commits. & \texttt{AtomicPublication} \\
\hline
\texttt{UnsafeReuseSerial} & Cancelled staged authority returns its consumed
    serial to the allocator. & \texttt{SerialsConsumed} \\
\hline
\tlbreak{UnsafeCancel}{Successor} & A late reservation cancellation alters a
    replacement namespace. & \texttt{ReplacementIntact} \\
\hline
\tlbreak{UnsafeRestore}{Successor} & Source rollback overwrites a replacement's
    unrelated payload. & \texttt{ReplacementIntact} \\
\hline
\tlbreak{UnsafeRefund}{Retained} & A detached metadata node loses its charge
    before disposal. & \texttt{RetainedNodesCharged} \\
\hline
\tlbreak{UnsafeReadiness}{Delivers} & A readiness hint makes hidden attachments
    visible. & \texttt{VisibleRequiresDelivery} \\
\hline
\tlbreak{UnsafeRepublish}{Observed} & Repeated reply observation revives a
    resource already closed by the caller. & \tlbreak{ObservedCloseStays}{Closed} \\
\hline
\tlbreak{UnsafeRollbackNeeds}{Admission} & Source restoration attempts fresh
    admission and strands escrow at the limit. &
    \tlbreak{RollbackRestores}{CapturedSource} \\
\hline
\tlbreak{Unsafe}{RemoteAbort} & A remote LP removes and reaps a context that may
    still execute. & \texttt{ReapOnlyOffCpu} \\
\hline
\tlbreak{UnsafeSpawn}{DuringAbort} & A sibling publishes after the domain-abort
    snapshot. & \texttt{AbortingThreadsDoomed} \\
\hline
\tlbreak{UnsafeCancelReleases}{Buffer} & Cancellation releases a mutable buffer
    before the operation becomes terminal. &
    \tlbreak{NonTerminalBuffer}{RemainsLoaned} \\
\hline
\tlbreak{UnsafeTimer}{Fire} & A watchdog reports timeout without rechecking a
    concurrently published work generation. & \texttt{TimeoutObservedNoWork} \\
\hline
\tlbreak{UnsafeRejected}{Release} & Wait admission rejection releases a pending
    operation's owner and loan. & \tlbreak{PendingRetains}{OwnerAndLoan} \\
\hline
\texttt{UnsafeReplay} & A timed-out accepted operation is submitted again. &
    \texttt{OneSubmission} \\
\hline
\tlbreak{UnsafeRestart}{Deadline} & A wake starts a fresh timeout budget. &
    \texttt{DeadlineUnchanged} \\
\hline
\tlbreak{UnsafeWakeAdopts}{Result} & A readiness hint adopts a result before
    confirmed completion. & \tlbreak{ResultNeeds}{TerminalProof} \\
\hline
\tlbreak{UnsafeParkWithout}{Recheck} & Completion precedes registration and the
    caller parks without checking it. & \texttt{ParkedTerminalHasWake} \\
\hline
\tlbreak{UnsafeAdmission}{Timeout} & Rejected admission is reported as timeout
    before the deadline. & \texttt{TimeoutRespectsDeadline} \\
\hline
\tlbreak{UnsafeObserve}{Join} & A delayed join attaches its old handle to the
    replacement occupying the same TID. &
    \tlbreak{ObserverMatches}{CapturedHandle} \\
\hline
\tlbreak{UnsafeSpawn}{Wrap} & Exhausted thread generation wraps and aliases a
    previously allocated identity. & \texttt{GenerationNeverReused} \\
\hline
\tlbreak{UnsafeGrowOver}{Budget} & Stack growth passes the admitted limit. &
    \texttt{CommittedWithinBudget} \\
\hline
\tlbreak{UnsafeKillWithout}{Release} & A dead domain retains frames that the
    stack-budget model requires it to release. & \texttt{DeadDomainsReleaseFrames} \\
\hline
\tlbreak{UnsafeBeginExclusive}{Map} & Exclusive DMA begins while CPU or other
    transfer authority still exists. &
    \tlbreak{ExclusiveDmaHasNo}{CpuAuthority} \\
\hline
\tlbreak{UnsafeClaim}{DetachedCell} & An available engine is used while its
    domain cell is missing. & \tlbreak{ClaimNeeds}{CompleteUnit} \\
\hline
\tlbreak{UnsafeReadiness}{Completes} & Unit readiness is treated as DMA
    completion. & \tlbreak{ReadinessIsNot}{Completion} \\
\hline
\tlbreak{UnsafeRejectDrops}{Preparation} & Rejected wait admission drops the
    original root, claim and pending pin. & \tlbreak{PreclaimRetains}{Preparation} \\
\hline
\texttt{UnsafeReplayMap} & A map error is used to authorize resubmission. &
    \texttt{MapSubmittedOnce} \\
\hline
\tlbreak{UnsafeRearmStarted}{Unmap} & A failure after backend claim rearms
    unmap. & \tlbreak{StartedFailureNot}{Rearmed} \\
\hline
\tlbreak{UnsafeDrop}{Resubmits} & Drop repeats an uncertain unmap attempt. &
    \tlbreak{UncertainAttemptNot}{Resubmitted} \\
\hline
\tlbreak{UnsafeForgetDma}{Confirmation} & Failed CPU unmap forgets confirmed
    DMA completion. & \texttt{ConfirmationSticky} \\
\hline
\tlbreak{UnsafeFailedShared}{Access} & Failed shared setup exposes CPU access. &
    \texttt{FailedSharedRestricted} \\
\hline
\tlbreak{UnsafeCpuUnmap}{DischargesDma} & CPU teardown releases uncertain DMA
    backing. & \tlbreak{PossibleDmaRetains}{Backing} \\
\hline
\tlbreak{UnsafeAbandon}{Releases} & Abandonment clears the complete owner and
    unit fence. & \tlbreak{AbandonmentRetains}{Owner} \\
\hline
\texttt{UnsafeWaitMasked} & A busy masked caller parks on readiness. &
    \tlbreak{WaitNeedsQualified}{Context} \\
\hline
\tlbreak{UnsafeFinalBusy}{Rearms} & A failed final public completion rearms
    unmap. & \tlbreak{PublicRejectionRetains}{Attempt} \\
\hline
\tlbreak{UnsafeCoherent}{Remap} & Fresh CPU mapping exposes Rust access while
    a coherent DMA pin remains. & \texttt{CpuOwnerNeedsNoDma} \\
\hline
\tlbreak{UnsafePublic}{Lookup} & Public lookup issues a policy-published
    connection. & \texttt{PublicIssuance} \\
\hline
\tlbreak{UnsafeDeferred}{Lookup} & A queued public lookup bypasses the
    publication class. & \texttt{PublicIssuance} \\
\hline
\tlbreak{UnsafePublic}{Replacement} & Public replacement displaces protected
    publication. & \texttt{ReplacementProtected} \\
\hline
\tlbreak{UnsafeMutationBefore}{Watch} & Policy changes before watch admission. &
    \texttt{BeforeWatchUnchanged} \\
\hline
\texttt{UnsafeHintRetires} & A readiness hint fabricates endpoint death. &
    \texttt{DeathNeedsExactResult} \\
\hline
\texttt{UnsafeStaleHint} & An old producer's result retires its replacement. &
    \texttt{DeathNeedsExactResult} \\
\hline
\tlbreak{UnsafeCloseSource}{First} & Source closes while its watch remains. &
    \texttt{WatchRetainsSource} \\
\hline
\tlbreak{UnsafeNewAt}{Capacity} & A new binding consumes replacement headroom. &
    \texttt{LiveBound} \\
\hline
\tlbreak{UnsafeAbort}{Disposes} & Abort disposes retained registry ownership. &
    \texttt{AbortRetainsOwners} \\
\hline
\tlbreak{UnsafeUnauthorized}{Retire} & Retirement skips its generation or actor
    check. & \texttt{RetirementAuthorized} \\
\hline
\tlbreak{UnsafeGeneration}{Wrap} & Exhaustion reuses an earlier generation. &
    \texttt{GenerationNotReused} \\
\hline
\tlbreak{UnsafePublish}{Exhausted} & Publication mutates policy beyond the
    signed NS generation range. & \tlbreak{PolicyGeneration}{Representable} \\
\hline
\tlbreak{UnsafeLookupDuring}{Publish} & Lookup runs between policy and registry
    publication. & \texttt{GrantGenerationCoherent} \\
\hline
\tlbreak{UnsafeUninstalled}{Publish} & A valid successor publishes before its
    installed-state commit. & \tlbreak{InstalledBefore}{Publication} \\
\hline
\texttt{UnsafeStaleCommit} & A second installer reuses a stale expectation. &
    \texttt{NoForkedInstallation} \\
\hline
\tlbreak{UnsafeEnrollment}{Fallback} & An installed-state read failure becomes
    enrollment. & \texttt{NoEnrollmentFallback} \\
\hline
\texttt{UnsafeOverwrite} & The same boot publishes a second policy. &
    \texttt{OnePublicationPerBoot} \\
\hline
\tlbreak{UnsafeIndependent}{Manifest} & Manifest and kernel gate select
    different policies. & \texttt{SharedPolicy} \\
\hline
\tlbreak{UnsafeWrong}{Recipient} & Preparation accepts the wrong private-key
    binding. & \texttt{RecipientBound} \\
\hline
\texttt{UnsafeAuthenticate} & A record bypasses authentication. &
    \texttt{VerifiedAuthenticated} \\
\hline
\tlbreak{UnsafeDevelopment}{Fallback} & Failed verification selects development
    trust. & \texttt{ExplicitDevelopment} \\
\hline
\tlbreak{UnsafeFailedCommit}{Publishes} & A failed commit grants publication
    authority. & \tlbreak{InstalledBefore}{Publication} \\
\hline
\tlbreak{UnsafeUnsupported}{Production} & Publication bypasses the unavailable
    protected boundary. & \texttt{ConditionalBoundary} \\
\hline
\texttt{UnsafeStaleClose} & A recycled numeric ASID is closed through an old
    generation handle. & \tlbreak{ReplacementSurvives}{StaleHandle} \\
\hline
\tlbreak{UnsafeDetachWith}{Lease} & A root detaches with a live operation lease. &
    \texttt{LeasesRetainRoot} \\
\hline
\tlbreak{UnsafeReleaseBefore}{Invalidation} & Backing is released without final
    invalidation acknowledgement. & \tlbreak{PhysicalRequires}{Invalidation} \\
\hline
\tlbreak{UnsafeRetry}{Physical} & A partially rejected physical walk is retried. &
    \texttt{PhysicalOnce} \\
\hline
\tlbreak{UnsafeAdopt}{Abandoned} & A controller reconstructs an abandoned owner. &
    \texttt{TerminalAbandonment} \\
\hline
\tlbreak{UnsafeRefund}{Partial} & Retained backing loses its original charge. &
    \texttt{BackingCharged} \\
\hline
\tlbreak{Unsafe}{Allocate} & A hardware ASID is reassigned before stale
    translations are invalidated. & \texttt{NoDirtyTagReuse} \\
\hline
\tlbreak{Drain}{Unsafe} & A deferred interrupt wake ignores its route generation.
    & \texttt{NoStaleWakeDelivery} \\
\hline
\tlbreak{UnsafeStale}{Unregister} & Cleanup from an old service generation removes
    its replacement. & \tlbreak{ReplacementSurvives}{StaleUnregister} \\
\hline
\tlbreak{UnsafeReplicate}{ToJoiner} & A non-selected peer supplies authority to a
    joining node. & \tlbreak{JoiningAcceptsOnly}{SelectedAnchor} \\
\hline
\tlbreak{UnsafeRestartForgets}{Admission} & Restart makes an incomplete joiner
    election-eligible. & \tlbreak{RestartPreserves}{Admission} \\
\hline
\tlbreak{UnsafeFinalize}{BeforeCatchUp} & Automatic FINALIZE bypasses proposed-member catch-up. &
    \tlbreak{ProposedMembers}{CaughtUp} \\
\hline
\tlbreak{UnsafeFinalize}{LearnerLeader} & A finalized learner retains leadership. &
    \tlbreak{LeaderIs}{Eligible} \\
\hline
\tlbreak{UnsafeFinalize}{OldOnly} & FINALIZE commits without the new voter majority. &
    \tlbreak{FinalizationHas}{BothQuorums} \\
\hline
\tlbreak{UnsafeReceive}{Lease} & The retired receive-time query gate admits an obsolete leader. &
    \tlbreak{ReadHasFence}{Proof} \\
\hline
\tlbreak{UnsafeRead}{Binding} & A query consumes another node/domain, term or membership episode. &
    \tlbreak{ReadBoundTo}{Admission} \\
\hline
\tlbreak{UnsafeReadBefore}{Commit} & A new leader reads before committing in its term. &
    \tlbreak{ReadHasCurrent}{Commit} \\
\hline
\tlbreak{UnsafeReadBefore}{Apply} & A read precedes application through the commit. &
    \tlbreak{ReadHasApplied}{Commit} \\
\hline
\tlbreak{UnsafeOldTerm}{Ack} & An old reply commits a replaced prefix. &
    \tlbreak{CommitHas}{CurrentPrefix} \\
\hline
\tlbreak{UnsafeRegress}{Success} & A delayed success discards confirmed progress. &
    \tlbreak{ConfirmedProgress}{Retained} \\
\hline
\tlbreak{Unsafe}{Backoff} & Rejection backs off below confirmed progress. &
    \tlbreak{BackoffRespects}{ConfirmedPrefix} \\
\hline
\tlbreak{UnsafeInvalid}{Credit} & An unqualified response gains credit. &
    \tlbreak{Acknowledgements}{HaveProof} \\
\hline
\tlbreak{UnsafeReject}{Application} & Rejected restore changes application state. &
    \tlbreak{RejectionPreserves}{Application} \\
\hline
\tlbreak{UnsafeReject}{Durable} & Rejected restore replaces durable bytes. &
    \tlbreak{RejectionPreserves}{Durable} \\
\hline
\tlbreak{UnsafeReject}{Progress} & Rejected restore advances progress. &
    \tlbreak{RejectionPreserves}{Progress} \\
\hline
\tlbreak{UnsafeReject}{Membership} & Rejected restore changes membership. &
    \tlbreak{RejectionPreserves}{Membership} \\
\hline
\tlbreak{UnsafeReject}{Join} & Rejected restore clears join admission. &
    \tlbreak{RejectionPreserves}{Join} \\
\hline
\tlbreak{UnsafeExpose}{Recovery} & Recovery exposes a rejected node. &
    \tlbreak{RecoveryExposure}{RequiresRestore} \\
\hline
\tlbreak{UnsafePublishBefore}{Restore} & Durable publication precedes restore. &
    \tlbreak{Publication}{RequiresRestore} \\
\hline
\tlbreak{UnsafeStartStale}{Flow} & A router admits a new SYN through an old
    snapshot after its admission lease expires. &
    \tlbreak{NewFlowsRequire}{FreshLease} \\
\hline
\tlbreak{UnsafeFallback}{ExistingPacket} & Snapshot-history eviction makes an
    established flow select a different backend. & \texttt{FlowBackendStable} \\
\hline
\tlbreak{UnsafeAdmitStale}{Readiness} & Readiness from an old deployment
    generation authorizes its replacement. &
    \tlbreak{CommittedPoliciesAre}{Derived} \\
\hline
\tlbreak{UnsafeEarlyPressure}{Move} & Placement pressure bypasses dwell or
    cooldown. & \texttt{PressureRespectsStability} \\
\hline
\texttt{UnsafeAdmit} & Placement reservations consume the protected reserve. &
    \tlbreak{ReservationsPreserve}{Reserve} \\
\hline
Flat session allocation & Retry and restart epochs alias one wire session. &
    \tlbreak{SessionIdentity}{Unique} \\
\hline
\texttt{AcceptDelayedSession} & An older delayed SYN regresses receive state. &
    \tlbreak{ReceiveSession}{Monotonic} \\
\hline
\texttt{UnsafeSetPolicy} & A non-administrator changes a policy rule. &
    \texttt{PolicyMutationAuthorized} \\
\hline
\tlbreak{UnsafeRedeemOther}{Principal} & A decision for one subject is redeemed
    for another. & \texttt{MintBoundToPrincipal} \\
\hline
\tlbreak{UnsafeRedeemStale}{Policy} & A decision survives a policy-version
    change. & \texttt{MintUsesCurrentPolicy} \\
\hline
\tlbreak{UnsafeRedeemStale}{Binding} & A decision for an old service generation
    reaches its replacement. & \texttt{MintTargetsCurrentBinding} \\
\hline
\tlbreak{UnsafeAmplify}{Rights} & Redemption mints rights absent from the
    decision. & \texttt{NoRightsAmplification} \\
\hline
\end{longtable}
\normalsize

The check suite expects all 107 negative configurations to fail for the named
reason. If one unexpectedly satisfies its invariant, the regression harness
has lost sensitivity and the suite fails.

\section{Concrete conformance and atomicity}

TLA+ actions are atomic. Rust code usually crosses several structures, locks,
or persistent writes. The conformance review must therefore identify a
linearization point or represent the concrete intermediate states explicitly.

The mapping uses two correspondence classes:

\begin{description}
    \item[Direct] The abstract action follows one concrete transition while
        omitting data irrelevant to the invariant. For example, a successful
        fenced unregister directly checks owner and generation.
    \item[Abstract] Several concrete operations are collapsed. This is valid
        only when intermediate states are unobservable or separately guarded.
        For example, scheduler retirement spans multiple tables, so the
        concrete \texttt{RETIREMENTS\_IN\_FLIGHT} state prevents address-space
        teardown during that interval.
\end{description}

Persistent ordering deserves the same care. In Raft admission, ``snapshot
membership and clear join fence'' is one model transition, but two durable
writes in Rust. Safety depends on snapshot flush preceding fence removal. A
crash between those writes must map to a modeled state that remains safe.

The conformance document is not yet a machine-checked refinement proof. It is
a reviewable refinement \emph{obligation}: it names the functions and state
that must be re-examined when either side changes. Code review should reject a
model update that merely hides a new concrete intermediate state.

\section{Validation}

\subsection{Running TLC}

The repository wrapper runs the complete bounded suite:

\begin{verbatim}
docs/tla/check.sh /path/to/tla2tools.jar
\end{verbatim}

Alternatively, \texttt{TLA2TOOLS\_JAR} may name the tools archive. The script:

\begin{itemize}
    \item runs all thirty-six safe configurations;
    \item enables action coverage and requires selected actions to have a
        positive successor count;
    \item rejects structural TLC warnings about variables, \texttt{EXCEPT}
        expressions, or incomplete successor states;
    \item runs 107 negative configurations and requires each one to fail on
        the named invariant after exercising the named unsafe action; and
    \item writes checkpoints and traces into a temporary directory so generated
        state does not become part of the source tree.
\end{itemize}

An invariant can pass because a transition is accidentally disabled. Required
action coverage detects this vacuity, although it does not prove that an action
is correct.

At the 6 September 2026 implementation synchronization point, representative
fast configurations produced the following complete state graphs:

\begin{longtable}{|p{0.32\textwidth}|r|r|r|}
\hline
\textbf{Model} & \textbf{Generated} & \textbf{Distinct} & \textbf{Depth} \\
\hline
\texttt{CharlotteIPC} & 6,118,645 & 1,254,153 & 13 \\
\hline
\texttt{CharlotteCQ} & 1,586,249 & 302,659 & 12 \\
\hline
\texttt{CharlotteTimedWait} & 22 & 15 & 8 \\
\hline
\texttt{CharlotteScheduler} & 4,308,577 & 819,972 & 22 \\
\hline
\texttt{CharlotteThreadJoin} & 46 & 39 & 10 \\
\hline
\tlbreak{CharlotteIpc}{Transaction} & 71 & 26 & 11 \\
\hline
\tlbreak{Charlotte}{ServiceLifecycle} & 689,321 & 74,232 & 37 \\
\hline
\texttt{CharlotteDMA} & 19,331,379 & 2,092,882 & 26 \\
\hline
\tlbreak{Charlotte}{Authorization} & 144,598 & 54,705 & 12 \\
\hline
\texttt{CharlotteRaftJoin} & 2,029 & 520 & 13 \\
\hline
\tlbreak{Charlotte}{ReliableMessage} & 30 & 21 & 5 \\
\hline
\tlbreak{CharlotteCluster}{Ingress} & 3,522,889 & 576,764 & 17 \\
\hline
\end{longtable}

State counts tell us how much of one configuration was explored. A
small model may encode the decisive identity argument, while a large model may
contain substantial independent bookkeeping. The invariant, abstraction, and
coverage matter more than graph size.

\subsection{Implementation validation}

TLC is followed by validation proportional to the affected code. The August
synchronization work, including the 15 August domain-abort and memory-ownership
repair, included:

\begin{itemize}
    \item host tests for \texttt{catten-graft}, including restart before durable
        membership and restart between membership persistence and fence clear;
    \item host tests for \texttt{charlotte-protocol-msg}, including disjoint
        retry/restart sessions and fail-closed exhaustion;
    \item host tests for ownership-aware runtime cleanup and the shared
        \texttt{charlotte-lifecycle} thread-join and timed-wait classifiers;
    \item AArch64 release builds and warning-denied Clippy for the changed
        \texttt{raft}, \texttt{relmsg}, and \texttt{ns} services;
    \item AArch64 and x86-64 kernel checks, plus warning-denied AArch64 Clippy,
        for the domain-abort publication fence and memory-ownership paths;
    \item a four-LP AArch64 QEMU boot whose authoritative result reported all
        eighteen registered self-tests passing, including owner-LP retirement,
        DMA/SMMU, completion cancellation, and service crash/restart; and
    \item formatting, shell syntax, and diff-whitespace checks.
\end{itemize}

Boot and multi-guest tests remain necessary when a change affects the concrete
scheduler, device, storage, or network integration. A TLC result cannot show
that QEMU delivers an interrupt, that an object-store flush reaches the
emulated disk, or that a manifest starts the intended services.

\section{What a successful check means}

A successful safe configuration means that TLC found no reachable violation
of the declared invariants within that specification and its finite constants.
It does \emph{not} by itself establish:

\begin{itemize}
    \item a proof that arbitrary larger bounds preserve the result;
    \item a mechanically checked refinement from Rust to TLA+;
    \item liveness, fairness, bounded latency, or eventual leader election
        unless explicitly expressed as temporal properties;
    \item correctness of omitted bytes, parsers, cryptography, memory mappings,
        hardware programming, or transport framing;
    \item correctness under storage failure modes below the modeled durable
        write interface; or
    \item freedom from ordinary implementation bugs such as arithmetic,
        aliasing, or unsafe-memory errors outside the modeled projection.
\end{itemize}

Conversely, a gap between code and model is actionable information. It can
mean one of three things:

\begin{enumerate}
    \item \textbf{Functionality is unsafe.} The Raft admission restart gap was
        of this kind and required a CharlotteOS change.
    \item \textbf{The model overclaims.} Atomic service replacement was of this
        kind; the model was corrected to admit temporary unavailability.
    \item \textbf{The behavior is outside scope.} It must be documented rather
        than silently attributed to the checked invariants.
\end{enumerate}

\begin{samepage}
\section{Workflow for changes}

Any change to ownership, generations, persistence, retry identity, lifecycle,
or distributed membership should be checked against this chapter's method.
A practical review checklist is:

\begin{enumerate}
    \item Search \repofile|docs/tla/CONFORMANCE.md| for the changed Rust
        function or state field.
    \item Decide whether the change adds an action, changes an enabling
        condition, moves a linearization point, or changes durable/volatile
        classification.
    \item Update the safe model and bounds. If repairing a discovered fault,
        retain the old behavior as an unsafe action and named negative
        configuration.
    \item Add the new action to the required-coverage list in
        \repofile|docs/tla/check.sh|.
    \item Run the entire TLC suite, including modules you did not edit, because shared
        documentation claims and layered Raft assumptions can drift together.
    \item Update the conformance map with exact Rust functions and ordering.
    \item Add implementation tests at the crash, retry, or stale-generation
        edges identified by the model.
    \item Run target validation and the relevant boot or multi-node suite.
    \item Record deliberate omissions and avoid converting a safety result into
        an unsupported liveness or availability claim.
\end{enumerate}
\end{samepage}

\endinput
